\section{Two Renderings Of One Constraint}

The stationary reader makes a demand before it supplies an equation. It
requires every admissible first residue to vanish. Classical notation often
writes this demand as
\[
  \delta A[\gamma;\eta] = 0
\]
for every permitted variation \(\eta\) of a path \(\gamma\). The notation is
helpful, but it is not the authority. The authority is the grammar that says
which variations are permitted, which endpoints remain fixed, and which
residue counts as first-order evidence of imbalance.

An endpoint is not held fixed because endpoints are sacred. It is held fixed
because the ledger has already committed it. If a variation could move the
committed boundary, the test would no longer be a test of the path between
the boundaries. It would be a new problem with a new record. The first
variation therefore reads only the interior degrees of freedom that the
gauge still permits to slip.

In the familiar variational calculation, the change of action separates
into an interior term and a boundary term. The boundary term vanishes when
the variation is zero at the endpoints. What remains is an integral or sum
against the arbitrary interior perturbation. If the first variation is zero
for every such perturbation, the coefficient of the perturbation must
vanish. The usual continuum expression is the Euler--Lagrange equation,
\[
  \frac{\partial L}{\partial q}
  - \frac{d}{dt}\frac{\partial L}{\partial \dot q} = 0 .
\]
The finite gauge reads the same obligation without requiring the continuum
to have been given first.

In the discrete reading, the path is a finite string of nodes. Each interior
node is touched from the left and from the right. A slip at that node changes
the cost of the left comparison and the cost of the right comparison. The
node is balanced exactly when those two changes cancel. The discrete
Euler--Lagrange equation is this nodewise balance. It says that no interior
node carries an unpaired first residual.

This is not a different theorem from stationarity. It is the same constraint
read at a different scale. The first-variation reading tests the whole path
against every admissible perturbation. The Euler--Lagrange reading inspects
the local balances that make all those tests vanish. If a node has a
nonzero residual, a variation supported near that node can read it. If every
supported variation reads zero, no node can retain a first residual. The
global vanishing and the local equations identify each other.

The grammar is doing important work here. It forbids a residue from hiding
behind a smooth phrase. A path may look continuous and still fail the local
balance. It also forbids the opposite mistake: a local equation may be
written down, but unless it is tied to the permitted variations and boundary
commitments, it has not yet been read as stationarity. The equation is not a
badge. It is the ledger entry left when the first variation has no remaining
debt.

The finite setting makes the equivalence especially sharp. Suppose an
interior node has a positive residual when pushed one way and a negative
residual when pushed the other. The path has not balanced that node. A
small supported variation exposes the imbalance immediately. Conversely, if
every interior node balances its incoming and outgoing costs, then every
finite supported variation decomposes into nodewise tests whose residues
sum to zero. The path is stationary because the grammar can find no first
residual anywhere it is allowed to look.

This also explains why the phrase ``least action'' must be handled with
care. A stationary path need not be globally least among every imaginable
path. The grammar has not licensed every imaginable path. It has licensed a
gauge, endpoints, and variations inside that gauge. The first claim is
therefore local and grammatical: no permitted first slip lowers or raises
the action to first order. Later positivity will say when this stationary
reading is also protected by a positive quadratic remainder.

The minimizing-variation example illustrates the distinction. A functional
with fixed endpoints is varied by a perturbation that vanishes at those
endpoints. Integration by parts moves the derivative off the perturbation
and leaves a term multiplying the free interior test. The classical
calculus says that this multiplier must vanish. The grammar says why that
move is legitimate: the boundary is pinned by commitment, the interior test
is arbitrary within the gauge, and a nonzero multiplier would be a readable
residue.

Thus the Euler--Lagrange equations are not imported as a law from outside
the finite grammar. They are the local reading forced by residual
vanishing. The finite gauge permits the test, carries the first variation,
identifies the local residuals, and forces their disappearance. Once that
has happened, the chapter can ask a more concrete question: what kind of
finite completion is sufficient to guarantee that no residual remains?
